% ============================================================
\section{Instance Solution Collection}

% ============================================================
\subsection{Overview}

The script \texttt{02\_get\_instance\_solutions.py} solves all training instances for a given problem class and records their optimal objective values. It must be run once before training, because the downstream learning pipeline requires ground-truth optimal values as supervision signals.

The script accepts three command-line arguments:

\begin{lstlisting}
python 02_get_instance_solutions.py <problem> [-j NJOBS] [-n NINST]
# problem: one of {setcover, cauctions, indset, mknapsack, ufacilities}
# njobs:   number of parallel worker threads (default 1)
# ninst:   maximum number of instances to solve (default 10000)
\end{lstlisting}

It scans the training directory for the selected problem, solves up to \texttt{ninst} instances in parallel using \texttt{njobs} threads, and writes a single JSON file \texttt{instance\_solutions.json} mapping each instance file path to its optimal objective value.

\subsubsection{Dependencies}

The script relies on two external libraries:
\begin{itemize}
    \item \textbf{Ecole} --- a Python library that wraps the SCIP solver and exposes a Gym-style reinforcement-learning interface for combinatorial optimization. The \texttt{ecole.environment.Configuring} environment is used here not for learning, but simply to solve each instance to optimality.
    \item \textbf{PySCIPOpt} --- accessed indirectly through Ecole via \texttt{model.as\_pyscipopt()}, used to retrieve the optimal objective value after solving.
\end{itemize}

% ============================================================
\subsection{Reward Function: \texttt{OptimalSol}}

Ecole environments accept a custom \emph{reward function} object that specifies what scalar feedback to return after each step. \texttt{OptimalSol} implements this interface:

\begin{lstlisting}
class OptimalSol:
    def before_reset(self, model): pass   # no setup needed

    def extract(self, model, done=False):
        pyscipopt_model = model.as_pyscipopt()
        if done:
            reward = pyscipopt_model.getObjVal(original=True)
        else:
            reward = None
        return reward
\end{lstlisting}

\begin{itemize}
    \item \texttt{before\_reset} is called by Ecole before each episode; it is a no-op here.
    \item \texttt{extract} is called after every \texttt{env.step()}. When \texttt{done=True} (i.e.\ the solver has finished), it returns the optimal objective value via \texttt{getObjVal(original=True)}. The flag \texttt{original=True} ensures the value is reported in the original, un-transformed objective space. Before the episode is done it returns \texttt{None}.
\end{itemize}

% ============================================================
\subsection{Worker Thread: \texttt{solve\_instance}}

\texttt{solve\_instance} is the function executed by each parallel worker thread. It loops over instances from a shared input queue until the queue is empty:

\begin{lstlisting}
def solve_instance(in_queue, out_queue):
    reward_fun = OptimalSol()
    while not in_queue.empty():
        instance = in_queue.get()
        env = ecole.environment.Configuring(
            scip_params={},
            observation_function=None,
            reward_function=reward_fun)
        env.reset(str(instance))         # load the .lp file
        _, _, solution, _, _ = env.step({})  # solve; solution = ObjVal
        out_queue.put({instance: solution})
\end{lstlisting}

\subsubsection{Ecole \texttt{Configuring} Environment}

The \texttt{Configuring} environment is the simplest Ecole environment: it loads an LP/MIP file, applies any SCIP parameter configuration supplied via \texttt{env.step(params\_dict)}, and solves to optimality. Passing an empty dict \texttt{\{\}} means ``apply no configuration changes, just solve.'' The five return values of \texttt{env.step()} are \texttt{(observation, action\_set, reward, done, info)}; only the third (\texttt{reward}) is used here, which equals the optimal objective value as defined by \texttt{OptimalSol.extract}.

\subsubsection{Thread Safety Note}

The \texttt{while not in\_queue.empty()} guard is not strictly atomic: two threads could both observe a non-empty queue before either calls \texttt{get()}. In practice this is harmless because \texttt{queue.Queue.get()} itself is thread-safe and raises \texttt{queue.Empty} if the queue is already drained; the worst outcome is that one thread exits the loop one iteration late.

% ============================================================
\subsection{Main Script Block}

\subsubsection{Instance Discovery}

After parsing arguments, the script selects the training directory for the chosen problem and collects all \texttt{.lp} files with \texttt{glob}:

\begin{lstlisting}
if args.problem == 'setcover':
    instance_dir = 'data/instances/setcover/train_400r_750c_0.05d'
# ... (one branch per problem type)
instances = glob.glob(instance_dir + '/*.lp')
num_inst = min(args.ninst, len(instances))
\end{lstlisting}

The directory names match exactly those produced by \texttt{01\_generate\_instances.py}.

\subsubsection{Producer--Consumer Architecture}

The script uses two \texttt{queue.Queue} objects and a pool of daemon threads:

\begin{enumerate}
    \item \textbf{Fill the input queue}: the first \texttt{num\_inst} instance paths are enqueued.
    \item \textbf{Launch workers}: \texttt{njobs} daemon threads are started, each running \texttt{solve\_instance(orders\_queue, answers\_queue)}.
    \item \textbf{Collect results}: the main thread blocks on \texttt{answers\_queue.get()} exactly \texttt{num\_inst} times, accumulating a \texttt{solutions} dict.
\end{enumerate}

\begin{lstlisting}
for instance in instances[:num_inst]:
    orders_queue.put(instance)

for i in range(args.njobs):
    p = threading.Thread(target=solve_instance,
                         args=(orders_queue, answers_queue),
                         daemon=True)
    p.start()

solutions = {}
for i in range(num_inst):
    solutions.update(answers_queue.get())
\end{lstlisting}

Because the threads are daemon threads, they are killed automatically if the main thread exits. The final \texttt{assert not p.is\_alive()} sanity-checks that all workers finished before the main thread reached that point (which is guaranteed because the main thread collected all \texttt{num\_inst} answers first).

\subsubsection{Output}

The \texttt{solutions} dict maps each instance file path (as a string) to its optimal objective value (a float). It is serialized as JSON:

\begin{lstlisting}
with open(instance_dir + "/instance_solutions.json", "w") as f:
    json.dump(solutions, f)
\end{lstlisting}

The file is written into the same training directory as the instances, e.g.:

\begin{lstlisting}
data/instances/setcover/train_400r_750c_0.05d/instance_solutions.json
\end{lstlisting}

% ============================================================
\subsection{Execution Flow Summary}

\begin{enumerate}
    \item Parse \texttt{problem}, \texttt{njobs}, \texttt{ninst}.
    \item Discover training \texttt{.lp} files; cap at \texttt{ninst}.
    \item Enqueue all instance paths into \texttt{orders\_queue}.
    \item Spawn \texttt{njobs} solver threads.
    \item Each thread: load instance $\to$ solve with SCIP via Ecole $\to$ record optimal value.
    \item Main thread collects all results and writes \texttt{instance\_solutions.json}.
\end{enumerate}
